매개변수 탐색(parametric search)는 최적해 문제(optimal problem)을 풀기 위한 강력한 도구로 이분 탐색을 이용해 최적해를 찾는 방법이다. 기본적으로 답을 탐색한다는 아이디어를 기반으로 하며, $O(\log X)$의 시간복잡도를 곱해 최적해 문제를 결정 문제(decision problem)으로 바꿀 수 있는 방법이다. 이번 장에서는 앞 장에서 배운 이분 탐색을 바탕으로 매개변수 탐색을 하는 방법을 배운다.

\section{답을 탐색한다는 아이디어}
이분 탐색이 어떤 값 $x$가 정렬된 배열에 있는지 빠르게 탐색하는 방법이었다. 이 과정에서 $x$와 \texttt{arr[mid]}를 비교하였다. 이를 일반적으로 확장한다면 어떤 조건 $f(x)$에 대해

\begin{itemize}
    \item $f(mid)$가 참일 때 답이 $mid$ 이상이고,
    \item $f(mid)$가 거짓일 때 답이 $mid$ 미만이면
\end{itemize}

항상 이분 탐색의 원리를 적용할 수 있다. 이러한 관찰에 착안하여 답을 탐색하겠다는 아이디어가 매개변수 탐색이다. 구체적인 내용은 아래 Observation을 참고하자.

\begin{obs}
    조건 $f(x)$를 '답이 $x$인 경우가 존재하는가?'로 잡자. 그러면 다음 조건을 만족할 때 답에 대한 이분탐색을 통해 최적해를 찾아낼 수 있다. 구체적인 부등식 조건은 가능한 값의 범위만 배수로 짧아짐이 보장되면 어떤 것이던 상관없다.

    \begin{itemize}
        \item 답이 $mid$인 경우가 존재할 때 최적해가 $mid$ 이상이고,
        \item 답이 $mid$인 경우가 존재하지 않을 때 최적해가 $mid$ 미만이면
    \end{itemize}
\end{obs}

이러한 방식을 매개변수 탐색이라고 한다. 매개변수 탐색을 시행하면 최적해를 찾는 문제를 답이 $x$일 수 있는지 판단하는 문제로 바뀐다. 이와 관련된 자세한 내용은 6.2절을 참고하라.

이제 예시 문제를 통해 구현 방법을 알아보자.

\begin{example} \label{ex: parametric-search}
    자연수로만 이루어진 어떤 배열 $A$가 주어질 때, $A_i$ 중 $x$보다 큰 수의 개수가 $k$개 이하가 되는 $x$의 최솟값을 구하여라.
\end{example}

물론 \Cref{ex: parametric-search}은 일반적인 이분 탐색을 통해서도 풀 수 있을 것이다. 하지만 연습을 위해 매개변수 탐색을 통해 풀어보자. 먼저 매개변수 탐색의 조건을 만족하는지 확인해야 한다. $x=mid$일 때 조건을 만족하면 답은 $mid$ 이하이다. 또한 $x=mid$일 때 조건을 만족하지 않는다면 답은 $mid$ 초과이다. 따라서 조건을 만족함을 확인할 수 있다. 그리고 이 부등식 조건을 그대로 코드에 구현하면 된다. \Cref{code: parametric-search}은 이를 바탕으로 구현한 코드이다. \Cref{code: parametric-search}에서 \texttt{isAnswer(x)}는 $x$가 답이 될 수 있는지 확인하는 함수이다.

\begin{code} \label{code: parametric-search}
\Cref{ex: parametric-search}의 정답 코드
\begin{lstlisting}
#include <bits/stdc++.h>
using namespace std;

int n, arr[101010], k;
const int X = 1e9;

bool isAnswer(int x) {
    int cnt = 0;
    for (int i = 1; i <= n; i++) if (arr[i] > x) cnt++;
    return cnt <= k;
}

int main() {
    cin >> n;
    for (int i = 1; i <= n; i++) cin >> arr[i];

    cin >> k;
    int st = 0, ed = X;
    while (st < ed) {
        int mid = (st + ed) / 2;
        if (isAnswer(mid)) ed = mid;
        else st = mid + 1;
    }

    cout << st << "\n";

    return 0;
}
\end{lstlisting}
\end{code}

\texttt{X}를 $A_i$의 범위라 하면 가능한 답은 $1$과 $X$ 사이이므로 \texttt{isAnswer} 함수가 호출되는 횟수는 $O(\log X)$번이고 전체 시간복잡도는 $O(N \log X)$가 된다.

\section{결정 문제와 최적해 문제}
PS의 문제는 크게 \textbf{결정 문제(decision problem)}와 \textbf{최적해 문제(optimal problem)}로 나눌 수 있다. 결정 문제는 어떤 조건을 만족하는 해가 존재하는지 판별하고 그 해를 출력하는 문제를 말하고, 최적해 문제는 어떤 조건을 만족하는 해들 중 가장 원하는 값이 큰 해를 찾아 출력하는 문제이다. 자명하게도 최적해 문제는 결정 문제보다 어렵다.

그렇다면 최적해 문제를 결정 문제로 환원할 수 있을까? 만약 어떤 최적해 문제의 답이 $k$ 이하일 때 항상 가능함이 보장된다면, $1$부터 차례대로 가능한지 확인해보면 되므로 결정 문제로 바꿀 수 있다. 즉, $k$번 결정 문제를 해결하여 최적해 문제의 답을 구할 수 있다.

위 방식을 개선시키기 위해 명제 $p_i$를 "답이 $i$일 때 가능하다"라고 정의해보자. 그러면 $i \leq k$이면 $p_i$는 참이고, $i > k$이면 $p_i$는 거짓이다. 즉, 우리가 $p_i$를 나열한 배열을 생각한다면 앞쪽은 전부 참이고, 뒤쪽은 전부 거짓이므로 \textbf{정렬된 배열}로 볼 수 있다(참인 경우 $1$을 적고, 거짓인 경우 $0$을 적는다고 생각하면 편하다).

따라서 이 배열 $p_i$에 이분 탐색을 적용할 수 있고, 이것이 매개변수 탐색이다. 즉, 배열의 크기를 $X$라고 하면 $O(\log X)$번의 결정 문제를 해결하여 최적화 문제의 답을 구할 수 있다. 또한 $X$는 가능한 답의 범위가 될 것이다.

매개변수 탐색은 $O(\log X)$의 시간을 곱해 최적해 문제를 결정 문제로 만들어주는 강력한 도구이다. 물론 $p_i$ 배열이 정렬된 상태, 즉 어떤 값을 기준으로 명제의 참/거짓이 뒤집혀야 한다.

\section{연습문제 A}
\begin{problem}
    \Cref{ex: parametric-search}을 이분 탐색으로 풀어라.
\end{problem}

\begin{problem}
    $N$개의 정수 $A_i$가 주어진다. 다음 값이 $K$ 이상이 되는 가장 큰 자연수 $x$를 출력하라. $N \leq 10^6$, $1 \leq A_i \leq 10^9$이다.
    $$\sum_{i=1}^N {\max(A_i - x, 0)}$$
\end{problem}

\begin{problem}
    $N$개의 정수 $A_i$가 주어진다. 다음 값이 $K$ 이상이 되는 가장 작은 자연수 $t$를 출력하라. $N \leq 10^6$, $1 \leq A_i \leq 10^9$이다.
    $$\sum_{i=1}^N {\left\lfloor{\frac{t}{A_i}}\right\rfloor}$$
\end{problem}

\section{연습문제 B}

\begin{problem}
    (1166번) 한 변의 길이가 $A$인 정육면체 $N$개를 각 변이 평행하도록 $L \times W \times H$ 크기의 직육면체 박스에 모두 넣으려고 한다. 가능한 $N, L, W, H$가 주어질 때 가능한 $A$의 최댓값을 구하여라. 절대 및 상대 오차는 $10^{-9}$까지 허용한다. $N, L, W, H$는 모두 $10^9$ 이하의 자연수이다.
\end{problem}

\begin{problem}
    (3090번) 정수 $N$개로 이루어진 수열 $A_i$에서 수열의 한 수를 골라 $1$만큼 감소시키는 연산을 $t$번 이내로 수행하여 인접한 수의 차이의 최댓값의 최솟값을 출력하라. 다른 말로,
    $$\max_{1 \leq i < N} {\left\vert{A_{i+1}-A_i}\right\vert}$$
    의 최솟값을 출력하라. $N \leq 10^5$, $t \leq 10^9$, $A_i \leq 10^9$이다.
\end{problem}

\begin{problem}
    (10227번, IOI10) $1$부터 $RC$까지의 자연수가 정확히 한 번씩 적킨 $R \times C$ 크기의 격자에서 홀수 $H$와 $W$가 주어질 때 $H \times W$ 크기의 부분 격자에 속한 원소의 중앙값의 최솟값을 출력하라. $R \leq 3000$, $C \leq 3000$이다.
\end{problem}

\begin{problem}
    (11233번) 세 수 $a, b, c$가 주어질 때 정삼각형 내부의 점 중 세 꼭짓점으로부터의 거리가 $a, b, c$인 점이 존재할 수 있는 정삼각형의 넓이를 출력하라. 만약 그러한 정삼각형이 없다면 $-1$을 출력하면 된다. \textbf{고정 소숫점 방식으로 소숫점 셋째 자리까지 출력해야 한다}. $a, b, c$는 모두 $100$ 이하의 양의 실수이다.
\end{problem}

\begin{problem}
    (15470번*) 중심이 같고 반지름이 다른 두 원이 있고, 반지름이 더 작은 원 위의 점 $N$개가 있다. 이 점들을 변이 교차하지 않도록 이은 다각형의 각 변의 길이 $l_i$와 반지름이 더 큰 원의 원주 길이 $c$가 주어질 때 바깥 원과 다각형 사이의 면적을 출력하라. $N \leq 10$, $10 \leq c \leq 1000$이다.
\end{problem}

\begin{problem}
    (18543번*) 두 명의 플레이어가 다음 규칙으로 게임을 한다.
    \begin{itemize}
        \item 선 플레이어가 처음 선택권을 갖는다.
        \item 선택권을 가진 플레이어가 $1$부터 $m$ 사이의 정수 중 하나를 뽑는다. 각 정수가 나올 확률은 모두 같다.
        \item 수를 뽑은 플레이어는 이 수를 상대에게 넘기고 자신의 선택권을 유지할지, 아니면 자신이 이 수를 갖고 상대에게 선택권을 넘길지 선택한다.
    \end{itemize}
    두 플레이어가 모두 최적으로 플레이할 때, $n$번의 숫자를 가진 이후 선 플레이어와 후 플레이어의 숫자 합 차이의 기댓값을 $d_n$이라 하자. 그러면 $m \geq 3$에서 $\lim_{n \rightarrow \infty} d_n$의 값이 유리수로 존재함이 보장된다. 여러분은 $m$이 주어질 때마다 $\lim_{n \rightarrow \infty} d_n$의 값을 $\textrm{mod}~10^9+7$에 대한 표현으로 출력해야 한다. $3 \leq m \leq 10^9$이고 $m$이 주어지는 횟수 $t$는 $5 \times 10^5$회 이하이다.
\end{problem}

\begin{problem}
    (20106번*, IOI13) $0$번부터 $N-1$번까지의 번호가 붙여져 있는 일렬로 나열된 $N$개의 문과 $N$개의 스위치가 있다. 각 스위치는 어느 문에 연결되어 있는 스위치인지 모른다. 또 스위치를 켰을 때 문이 열려야 정상이지만, 일부 스위치의 경우에는 켰을 때 문이 닫히는 오류가 있다. 여러분은 최대 $70000$번 스위치를 조정한 후 처음으로 닫힌 문이 몇 번 문인지 확인할 수 있다. 이 결과를 토대로 $i$번쨰 스위치가 문이 열리기 위해 켜져야 하는지(1) 꺼져야 하는지(0) 출력하고 $i$번 스위치와 연결된 문의 번호도 출력해야 한다.
    \begin{itemize}
        \item (서브테스크 1) $N \leq 5000$, $i$번 스위치는 $i$번 문과 연결되어 있다.
        \item (서브테스크 2) $N \leq 5000$, 모든 스위치는 꺼졌을 때 문이 열린다.
        \item (서브테스크 3) $N \leq 100$
        \item (서브테스크 4) $N \leq 2000$
        \item (서브테스크 5) $N \leq 5000$
    \end{itemize}
\end{problem}

\begin{problem}
    (23791번) 길이 $N$의 두 배열 $A$와 $B$가 주어질 때, 다음 쿼리 $Q$개를 처리하는 프로그램을 작성하라. $N \leq 10^5$, $Q \leq 10^5$이고 $1 \leq A_i \leq 2^{31}-1$, $1 \leq B_i \leq 2^{31}-1$이다.
    \begin{itemize}
        \item $i, j, k$를 입력받아 $A_1$부터 $A_i$, 그리고 $B_1$부터 $B_j$ 까지의 숫자들 중 $k$번째로 작은 수를 출력하라. 다른 말로 $A_1$부터 $A_i$, 그리고 $B_1$부터 $B_j$ 까지의 숫자들 중에서 $x$ 이하인 수의 개수가 $k$인 가장 작은 $x$를 출력하라.
    \end{itemize}
\end{problem}

\begin{problem}
    (23107번*) 주어진 길이 $N$의 수열 $a_i, b_i, c_i$에 대해 다음 조건을 만족하는 순서쌍 $S=(s_0, s_1, \cdots, s_r)$을 모두 생각하자.
    \begin{itemize}
        \item $r$은 자연수
        \item $s_0 = 0, s_r = N$
        \item $s_i < s_{i+1}$ $(0 \leq i < r)$
    \end{itemize}
    $f(S)$를 다음과 같이 정의할 때 가능한 모든 $S$에 대해 $f(S)$의 최댓값을 출력하라. $N \leq 2 \cdot 10^5$이고 $a_i, b_i, c_i$는 절댓값이 $10^9$ 이하인 정수이다.
    $$f(S) = \min_{0 \leq i < r} {\left \{ b_{s_i} + c_{s_{i+1} - 1} + \sum_{s_i \leq j < s_{i+1}} {a_j} \right \}}$$
\end{problem}

\begin{problem}
    (27811번*) 양쪽 더미에 $L$개와 $R$개의 돌이 있다. $i$번째 라운드에 더 많은 개수의 돌이 있는 더미에서 $i$개의 돌을 가져간다. 만약 돌의 개수가 같으면 왼쪽 더미에서 $i$개의 돌을 가져간다. 만약 $i$번쨰 라운드에 양쪽 더미 모두에 있는 돌의 개수가 $i$개 미만이 되면 돌을 가져가지 못하므로 $i$번째 라운드가 진행되지 못하고 게임이 종료된다. 진행된 라운드의 숫자와, 게임 종료 후 양쪽 더미에 남은 돌의 개수를 출력하라. $T$번 $L$과 $R$이 주어졌을 때 출력해야 한다. $T \leq 1000$, $L \leq 10^{18}$, $R \leq 10^{18}$이다.
\end{problem}

\section{연습문제 C}
\begin{problem}
    (1557번) $1$이 아닌 제곱수로 나누어지지 않는 자연수를 크기 순으로 나열할 때 $K$번째로 나타나는 수를 출력하라. $K \leq 10^9$이다.
\end{problem}

\begin{problem}
    (2814번) 소수 $p$에 대해 가장 작은 소인수가 $p$인 자연수를 크기 순으로 나열할 때 $N$번째로 나타나는 수를 출력하라. $N \leq 10^9$, $p \leq 10^9$이다.
\end{problem}

\begin{problem}
    (5563번) 주어진 $N$개의 점을 두 그룹으로 분할하려고 한다. 각 그룹에는 한 개 이상의 점이 포함되어야 한다. 같은 그룹에 속하는 임의의 두 점 사이의 거리의 최댓값의 최솟값을 출력하라. 두 점 $(x_1, y_1)$과 $(x_2, y_2)$ 사이의 거리는 $|x_1 - x_2| + |y_1 - y_2|$로 계산된다. $3 \leq N \leq 10^5$이고 점의 $x$좌표와 $y$좌표는 절댓값이 $10^5$ 이하인 정수이다.
\end{problem}

\begin{problem}
    (9014번) 주어진 $N$개의 점을 두 개의 기준점 $(a_1, b_1)$과 $(a_2, b_2)$ 중 적어도 하나와 연결하려고 한다. 그런데 각각의 기준점은 적어도 $K$개의 점과 연결되어야 한다. 여러분은 $N$개의 점과 기준점과의 거리의 최댓값을 최소화해야 한다. 가능한 최솟값을 소숫점 첫째 자리에서 반올림한 값을 출력하라. 두 점 $(x_1, y_1)$과 $(x_2, y_2)$ 사이의 거리는 $|x_1 - x_2| + |y_1 - y_2|$로 계산된다. $2 \leq N \leq 10^5$, $K \leq N/2$이고 점의 $x$좌표와 $y$좌표는 절댓값이 $10^5$ 이하인 짝수이다.
\end{problem}

\begin{problem}
    (16655번) $N$명의 사람이 $S$달러를 나누어 가지려고 한다. 하지만 $N$명의 사람 각각이 요구하는 금액 $a_i$의 합이 $S$보다 커 문제가 되고 있다. 다음과 같이 \textit{공정한 분할}과 \textit{올바른 분할}을 정의할 때 \textit{공정한 분할}의 결과를 출력하라. $N \leq 5000$, $a_i \leq 5000$, $S \leq \sum_{i=1}^N {a_i}$이다.
    \begin{itemize}
        \item \textit{공정한 분할}은 합이 $S$인 수열 $c_i$들 중 모든 $(i, j)$ 쌍에 대해 $(a_i, a_j)$에 대한 합 $c_i + c_j$의 \textit{올바른 분할}이 $(c_i, c_j)$인 수열이다.
        \item $(x, y)$에 대한 합 $z$의 \textit{올바른 분할}은 $x+y \geq z$인 경우에 아래와 같이 정의된다.
        $$\begin{cases} 
            (\frac{z}{2}, \frac{z}{2}) & x \geq z, y \geq z \\
            (\frac{x}{2}, z - \frac{x}{2}) & x \leq z \leq y \\
            (z - \frac{y}{2}, \frac{y}{2}) & y \leq z \leq x \\
            (\frac{x+z-y}{2}, \frac{y+z-x}{2}) & x \leq z, y \leq z
        \end{cases}$$
    \end{itemize}
\end{problem}

\begin{problem}
    (18592번**) 중심이 $(x_i, y_i, z_i)$이고 반지름이 $r_i$인 $N$개의 구를 생각하자. $N(N-1)/2$개의 $(i, j)$ 쌍에 대하여 $d(i, j)$를 두 구 사이의 최단거리로 정의하자(겹치면 $0$이다). 수식으로는 아래와 같이 나타낼 수 있다.
    $$d(i, j) = \max \left( 0, \sqrt{(x_i - x_j)^2 + (y_i - y_j)^2 + (z_i - z_j)^2} - r_i - r_j \right)$$
    모든 $(i, j)$ 쌍에 대한 $\lceil d(i, j) \rceil$를 크기 순으로 나열할 때 가장 작은 $k$개를 증가하는 순서로 출력하라. $N \leq 10^5$, $k \leq 300$, $0 \leq x_i, y_i, z_i\leq 10^6$, $1 \leq r_i \leq 10^6$이다.
\end{problem}

\begin{problem}
    (26610번) $N$개의 돌이 있는 더미에서 $i$번째 턴에 $1 \leq k \leq T$인 자연수 $k$에 대해 $(X+i) \times k$개의 돌을 가져가야 하는 게임을 생각하자. 만약 더 이상 조건을 만족하게 돌을 가져갈 수 없으면 게임이 종료된다. 여러분은 다음 쿼리 $Q$개에 대해 답해야 한다. $Q \leq 3 \times 10^5$, $1 \leq X < N \leq 5 \times 10^8$, $2 \leq T \leq 10^8$이다.
    \begin{itemize}
        \item 이 게임이 종료되었을 때 가장 적은 개수의 돌을 남겨야 하고 만약 그러한 방법이 여러 개라면 가장 빠르게 게임을 끝내야 한다. 최적의 방식으로 플레이했을 때 돌을 가져가는 횟수를 출력하라.
    \end{itemize}
\end{problem}